
En este apartado abordaremos la parte correspondiente a la práctica 2 de la asignatura.
Se puede observar en la figura \ref{fig:dfd_publi} el Diagrama de Flujo de Datos (DFD) 
de nivel 1 del subsistema de publicidad. En este diagrama se representan los diferentes 
procesos, almacenes de datos, agentes externos y flujos de datos que componen el subsistema.

\begin{figure}[h!]
    \centering
    % Ajusta el tamaño del diagrama al ancho del documento (opcional)
    \shorthandoff{<>}
    \resizebox{\textwidth}{!}{%
    \begin{tikzpicture}[node distance=5cm, auto]

        % ======================================================
        % === 1. AGENTES EXTERNOS ===============================
        % ======================================================
        % Los "agentes" representan entidades externas al sistema
        % Ejemplo: usuarios, sistemas externos, administradores, etc.

        \node[agente, minimum width=4cm, minimum height=1.2cm] (Agente1) {Usuario};
        %\node[agente, right of=Agente1, node distance=14cm] (Agente2) {Sistema Externo};

        % ======================================================
        % === 2. ALMACÉN DE DATOS ===============================
        % ======================================================
        % Representan bases de datos o repositorios.
        % Suelen ser cilindros o rectángulos especiales.

        \node[almacen, below of=Agente1, node distance=10cm, xshift=7cm, minimum width=6cm] (BD1) {BD EKIS};

        % ======================================================
        % === 3. PROCESOS (FUNCIONES O RFs) =====================
        % ======================================================
        % Cada proceso es una función, requisito funcional (RFx.x) o tarea.
        % Se ubican entre los agentes y los datos.

        \node[proceso, below of=Agente1, node distance=3cm, xshift=-3cm, minimum width=4.5cm] (P1) {RF1.1: Crear publicación};
        \node[proceso, below of=P1, node distance=4cm, minimum width=4.5cm] (P2) {RF1.2: Modificar publicación};
        \node[proceso, right of=P1, node distance=8cm, minimum width=4.5cm] (P4) {RF1.4: Dar me gusta a una publicación};
        \node[proceso, right of=Agente1, node distance=6cm, xshift=3cm, minimum width=4.5cm] (P5) {RF1.5: Eliminar publicación};
        \node[proceso, below of=P4, node distance=2cm, minimum width=4.5cm] (P3) {RF1.3: Listar publicaciones};

        % ======================================================
        % === 4. FLUJOS DE DATOS ================================
        % ======================================================
        % Usa \draw[flujo] para representar conexiones.
        % Puedes añadir etiquetas con node[pos=..., ...]{Texto}.
        % Ejemplos:
        %   RDE = Requisito de entrada (data entrada)
        %   RDS = Requisito de salida (data salida)
        %   RDW = Escritura a BD
        %   RDR = Lectura desde BD

        % ------------------------------------------------------
        % Ejemplo de flujo entre Agente y Proceso (RDE)
        % ------------------------------------------------------
        \draw[flujo, ->] (Agente1.south west) ++(0.2,0) |- ++(0,-1) -| node[pos=0.7, right] {RDE 1.1} (P1.north);
        \draw[flujo, ->] (Agente1.west) |- ++(0,-1) -| node[pos=0.53] {RDE 1.2} (P2.north west);
        %RF1.3 no tiene RDE
        \draw[flujo, ->] (Agente1.south east) -| node[pos=0.6, left] {RDE 1.3} (P4.north);
        \draw[flujo, ->] (Agente1.east) -- node[pos=0.4, above] {RDE 1.4} (P5.west);
        
        % Proceso a Agente (RDS)
        %RF1.1 no tiene RDS
        \draw[flujo, ->] (P2.west) -- ++(-0.8,0) -| ++(-0.8,7) -- node[pos=0.4, above] {RDS 1.1} (Agente1.west);
        \draw[flujo, ->] (P3.west) -- node[pos=0.55, below] {RDS 1.2} ++(-1.5,0) |- ++(-0.97,2.5) -- (Agente1.south);
        \coordinate (elbow) at ($(P4.north west)+(0,1.875)$); 
        \draw[flujo, ->] (P4.north west) -- (elbow) node[pos=0.2, right] {RDS 1.3};

        % ------------------------------------------------------
        % Ejemplo de flujo entre Proceso y Almacén de Datos (RDW)
        % ------------------------------------------------------
        \draw[flujo, ->] (P1.east) -- ++(1.0,0) |- ++(1.0,-7) node[pos = 0.465] {RDW 1.1} -- (BD1.west);
        \draw[flujo, ->] (P2.south east) -- node[pos = 0.15, left] {RDW 1.2} ++(0, -3) -- (BD1.south west);
        %RF1.3 no genera ningún RDW
        \draw[flujo, ->] (P4.east) |- ++(-1.6, -3) -- node[pos = 0.1] {RDW 1.3} (BD1.north);
        \draw[flujo, ->] (P5.east) -- ++(0.5, 0) |- node[pos=0.68, below] {RDW 1.4} (BD1.east);
        %\draw[flujo, ->] (P5) -- (BD1); %ESTA LINEA

        % ------------------------------------------------------
        % Ejemplo de flujo desde Almacén a Proceso (RDR)
        % ------------------------------------------------------
        %RF1.1 no genera ningún RDR
        \draw[flujo, ->] (BD1.west) |- ++(-2,-1) -| node[pos=0.6, left] {RDR 1.1} (P2.south);
        \draw[flujo, ->] (BD1.north west) ++(0,-0.2) -- ++(0,2) node[pos=0.8, right] {RDR 1.2} -| (P3.south);
        \draw[flujo, ->] (BD1.north east) -- node[pos=0.2, left] {RDR 1.3} ++(0,6.49) -- (P4.east);

        % ------------------------------------------------------
        % Ejemplo de flujo entre Procesos
        % ------------------------------------------------------
        %\draw[flujo, ->] (P1.south) -- node[midway, right] {RDF interno} (P2.north);

        % ------------------------------------------------------
        % Ejemplo de flujo entre Proceso y Agente externo
        % ------------------------------------------------------
        %\draw[flujo, ->] (P4.east) -- ++(2,0) |- node[pos=0.3, right] {RDS 1.4} (Agente2.south);

    \end{tikzpicture}
    } % Fin del resizebox
    \caption{DFD1 — Subsistema de publicaciones}
    \label{fig:dfd_publi}
\end{figure}
